\chapter*{Synthesis}
\addcontentsline{toc}{chapter}{Synthesis}
\label{chap:summary_english}

The valuation of financial liabilities under the Solvency II framework requires heavy stochastic projections, generating calculation times that are often incompatible with the operational constraints of insurance companies. This study aims to resolve this dimensionality issue by exploring various aggregation techniques applied to a life insurance liability portfolio. The objective is to aggregate policy information while preserving the accuracy of regulatory risk indicators, specifically the \textit{Best Estimate} (BE) and the Solvency Capital Requirement (SCR).

Faced with the massive volume of life insurance portfolios, the standard aggregation method relies on the creation of \textit{Model Points} (MP). These aggregated portfolios group policies with similar risk profiles. The major technical challenge lies in the calibration of this compression: excessive aggregation leads to a smoothing of guarantees, altering the calculation of the BE and SCR. To align with market operational standards, this study imposes a strict compression constraint, limiting the portfolio size to a volume between 800 and 1,200 \textit{Model Points}.

To conduct this analysis, a synthetic reference portfolio of 1,000,000 policies was generated, faithfully reflecting the structure of the French market. This stochastic generator is based on a generation-based approach, ensuring temporal and financial consistency of the simulated technical characteristics. The statistical distributions used were calibrated on empirical public data:
\begin{itemize}
    \item \textbf{Policyholder Age:} Beta distribution, reproducing the demographic asymmetry of life insurance holders;
    \item \textbf{Gender:} Bernoulli distribution, parameterized according to the observed sector distribution;
    \item \textbf{Policy Seniority:} Gaussian Mixture Model (GMM), selected to capture different historical waves of inflows;
    \item \textbf{Mathematical Provisions (MP):} Log-Normal distribution, capturing the asymmetry of wealth distribution and the concentration of assets in high-capital contracts;
    \item \textbf{Conditional Technical Characteristics:} The minimum guaranteed rate (MGR) and the share of unit-linked (UL) funds are conditioned on the year of subscription to reflect changes in the economic environment.
\end{itemize}

In order to compare different dimension reduction logics, a strict aggregation based on identical characteristics serves as the \textit{Baseline}. Inducing no loss of information, it is compared to five distinct approaches:
\begin{itemize}
    \item \textbf{Banding:} A standard deterministic method segmenting the portfolio according to predefined grids applied to key variables (age, seniority, MGR);
    \item \textbf{K-Means Algorithm:} An unsupervised learning method using spatial partitioning to minimize the intra-class variance of static attributes;
    \item \textbf{DBSCAN and HDBSCAN Algorithms:} Spatial approaches based on local density, allowing for the isolation of atypical profiles through noise management. Processing discrete variables requires the prior implementation of uniform white noise (\textit{Dithering});
    \item \textbf{Cash-Flow Matching:} A dynamic aggregation method applied not to static characteristics at the closing date, but to the morphology of future cash flow run-offs. The \textit{K-Means} algorithm is used to group policies with similar flow trajectories;
    \item \textbf{A Posteriori Calibration:} An approach consisting of grouping policies according to their financial characteristics. Surrender laws are then calibrated for each group to reproduce the reference portfolio's BE identically in a deterministic liability-only projection.
\end{itemize}

The hyperparameter selection protocol is based on an exhaustive search (\textit{Grid Search}), aiming to identify the Pareto optimum for each algorithm. The goal is to achieve the best compromise between the imposed operational compression constraint and the minimization of the error on the \textit{Best Estimate} (BE), initially evaluated using a deterministic liability-only model.

Table \ref{tab:synthese_calibration_finale_synthese_english} presents the results obtained on the liability-only model. It shows that all grouping methods prove to be very accurate in this baseline framework, with negligible errors compared to the initial full portfolio.

\begin{table}[H]
    \centering
    \renewcommand{\arraystretch}{1.4}
    \begin{tabular}{|l|r|r|r|}
        \hline
        \textbf{Aggregation Method} & \textbf{Nb MP} & \textbf{Liability-Only BE} & \textbf{Liability-Only Error} \\
        \hline
        \textit{Baseline} (Reference) & 26,272 & €263.95 Bn & \textit{Ref.} \\
        Dynamic \textit{Banding} & 1,172 & €264.03 Bn & $+0.0292\%$ \\
        Robust \textit{K-Means} & 1,039 & €263.96 Bn & $+0.0035\%$ \\
        \textit{DBSCAN} & 1,047 & €263.98 Bn & $+0.0124\%$ \\
        \textit{HDBSCAN} & 967 & €264.06 Bn & $+0.0398\%$ \\
        \textit{Cash-Flow Matching} & 1,000 & €263.94 Bn & $-0.0049\%$ \\
        \textit{A Posteriori} Calibration & 222 & €263.95 Bn & $0.0000\%$ \\
        \hline
    \end{tabular}
    \caption{Deterministic liability-only evaluation and associated aggregation errors.}
    \label{tab:synthese_calibration_finale_synthese_english}
\end{table}

The strong compression of the volume, divided by more than twenty-five compared to the reference portfolio, occurs without significant deformation of the BE, with the maximum error peaking at $+0.0398\%$. Furthermore, the perfect restitution of the \textit{a posteriori} calibration stems mechanically from its construction, as the surrender laws are specifically adjusted to replicate the reference deterministic commitment.

Before testing the robustness of these models against regulatory shocks, the analysis of execution times (Table \ref{tab:runtimes_synthese_english}) quantifies the performance gain obtained during the transposition to a full stochastic ALM environment.

\begin{table}[H]
    \centering
    \renewcommand{\arraystretch}{1.4}
    \begin{tabular}{|l|r|r|r|}
        \hline
        \textbf{Aggregation Method} & \textbf{Volume (MP)} & \textbf{Execution Time} & \textbf{Reduction Factor} \\
        \hline
        \textit{Baseline} (Reference) & 26,272 & $\approx$ 8 h 00 (28,783 s) & \textit{Ref.} \\
        \textit{HDBSCAN} & 967 & $\approx$ 10 min 00 (600 s) & $\div 48$ \\
        \textit{Cash-Flow Matching} & 1,000 & $\approx$ 10 min 27 (627 s) & $\div 46$ \\
        \textit{K-Means} Robust & 1,039 & $\approx$ 11 min 11 (671 s) & $\div 43$ \\
        \textit{Banding} Dynamique & 1,172 & $\approx$ 11 min 51 (711 s) & $\div 40$ \\
        \textit{A Posteriori} Calibration & 222 & $\approx$ 2 min 38 (158 s) & $\div 182$ \\
        \hline
    \end{tabular}
    \caption{Compared calculation times in stochastic ALM projection (5,000 scenarios).}
    \label{tab:runtimes_synthese_english}
\end{table}

The reduction in volume divides processing times by a factor of 46 at 1,000 MP, bringing an exhaustive eight-hour projection down to about ten minutes. This performance gain allows for the multiplication of sensitivity calculations on standard infrastructures.

Validating the aggregation requires analyzing the balance sheet under the central stochastic scenario. The ALM model integrates non-linear interactions between financial returns and policyholder behavior through profit-sharing mechanics. Table \ref{tab:resultats_central_alm_english} shows the BE and SCR evaluation within this study framework.

\begin{table}[H]
    \centering
    \renewcommand{\arraystretch}{1.4}
    \begin{tabular}{|l|r|r|r|r|}
        \hline
        \textbf{Aggregation Method} & \textbf{Central BE} & \textbf{BE Error} & \textbf{Global SCR} & \textbf{SCR Error} \\
        \hline
        \textit{Baseline} (Reference) & €209.93 Bn & \textit{Ref.} & €26.91 Bn & \textit{Ref.} \\
        Dynamic \textit{Banding} & €209.02 Bn & $-0.43\%$ & €27.47 Bn & $+2.10\%$ \\
        \textit{HDBSCAN} & €209.05 Bn & $-0.42\%$ & €27.47 Bn & $+2.10\%$ \\
        Robust \textit{K-Means} & €208.95 Bn & $-0.46\%$ & €27.49 Bn & $+2.17\%$ \\
        \textit{Cash-Flow Matching} & €208.41 Bn & $-0.72\%$ & €27.70 Bn & $+2.94\%$ \\
        \textit{A Posteriori} Calibration & €198.31 Bn & $-5.53\%$ & €19.27 Bn & $-28.38\%$ \\
        \hline
    \end{tabular}
    \caption{Global evaluation of liabilities and SCR in stochastic ALM projection.}
    \label{tab:resultats_central_alm_english}
\end{table}

Spatial algorithms (\textit{HDBSCAN}, \textit{K-Means}) and \textit{Banding} demonstrate high precision, limiting the global SCR deviation below the $2.20\%$ threshold. To isolate the origin of aggregation biases, the error decomposition across the main risk modules is presented in Table \ref{tab:resultats_modules_english}.

\begin{table}[H]
    \centering
    \renewcommand{\arraystretch}{1.4}
    \begin{tabular}{|l|r|r|r|}
        \hline
        \textbf{Aggregation Method} & \textbf{Market SCR Error} & \textbf{Life SCR Error} & \textbf{Global SCR Error} \\
        \hline
        Dynamic \textit{Banding} & $+1.89\%$ & $+3.55\%$ & $+2.10\%$ \\
        \textit{HDBSCAN} & $+1.87\%$ & $+3.71\%$ & $+2.10\%$ \\
        Robust \textit{K-Means} & $+1.92\%$ & $+3.87\%$ & $+2.17\%$ \\
        \textit{Cash-Flow Matching} & $+2.42\%$ & $+6.50\%$ & $+2.94\%$ \\
        \textit{A Posteriori} Calibration & $-52.82\%$ & $+83.88\%$ & $-28.38\%$ \\
        \hline
    \end{tabular}
    \caption{Error decomposition across the main risk modules.}
    \label{tab:resultats_modules_english}
\end{table}

The dynamic approach of \textit{Cash-Flow Matching} shows a marked overestimation of the Life SCR ($+6.50\%$). Based exclusively on the central morphology of flows, this method obscures the heterogeneity of tax seniority within the groups, generating a mechanical over-reaction to structural lapse assumptions. Finally, the \textit{a posteriori} calibration method proves ineffective in a stochastic environment because fixed behavioral laws fail to replicate the reality of the portfolio in a stochastic environment.

Evaluating model resilience requires applying sensitivities to the economic environment. The portfolio is thus subjected to two parallel shifts of the risk-free rate curve published by EIOPA, amounting to $-100$ and $+100$ basis points respectively. This step tests the resistance of grouping methods against different economic environments (Table \ref{tab:sensibilites_scr_synthese_english}).

\begin{table}[H]
    \centering
    \renewcommand{\arraystretch}{1.4}
    \resizebox{\textwidth}{!}{
    \begin{tabular}{|l|r|r|r|r|}
        \hline
        \multirow{2}{*}{\textbf{Aggregation Method}} & \multicolumn{2}{c|}{\textbf{Decrease (-100 bps)}} & \multicolumn{2}{c|}{\textbf{Increase (+100 bps)}} \\
        \cline{2-5}
        & \textbf{Lapse SCR Error} & \textbf{Global SCR Error} & \textbf{Lapse SCR Error} & \textbf{Global SCR Error} \\
        \hline
        Dynamic \textit{Banding} & $+7.55\%$ & $+1.17\%$ & $+4.55\%$ & $+1.23\%$ \\
        \textit{HDBSCAN} & $+8.22\%$ & $+1.16\%$ & $+4.68\%$ & $+1.24\%$ \\
        Robust \textit{K-Means} & $+11.56\%$ & $+1.21\%$ & $+4.72\%$ & $+1.25\%$ \\
        \textit{Cash-Flow Matching} & $+1.11\%$ & $+1.39\%$ & $+6.23\%$ & $+1.27\%$ \\
        \textit{A Posteriori} Calibration & $+1779.84\%$ & $-19.19\%$ & $+0.47\%$ & $-37.16\%$ \\
        \hline
    \end{tabular}
    }
    \caption{Relative errors on Global SCR and Lapse SCR according to interest rate curve shifts.}
    \label{tab:sensibilites_scr_synthese_english}
\end{table}

The capital requirement for lapse risk shows strong asymmetry depending on the financial environment. In a low-rate scenario, asset returns are absorbed by the cost of minimum guaranteed rates (MGR), so future portfolio profitability is limited. Applying a regulatory lapse shock (upward shock or mass lapse) does not decrease the margin much and therefore does not significantly affect the SCR. Conversely, an increase in interest rates restores a very positive return differential. The lapse shock then prevents a significant future margin. It is the materialization of this financial opportunity cost that justifies the sharp increase in the upward Lapse SCR.

Spatial algorithms, by preserving the granularity of guaranteed rates and static variables driving structural lapse laws (such as seniority), capture this loss of value with an error contained around $+4.70\%$ in the upward scenario. In contrast, the \textit{Cash-Flow Matching} method shows a higher deviation ($+6.23\%$). By smoothing individual characteristics in favor of an aggregation based solely on the morphology of central flows, this approach alters the application of base structural lapse rates, generating a portfolio over-reaction to regulatory shocks.

In conclusion, this study validates the use of unsupervised learning algorithms for aggregation within ALM modeling, while providing critical insight into industry practices:

\begin{itemize}
    \item \textbf{Robustness of the empirical approach:} The \textit{Banding} method demonstrates marked resilience to asymmetric shocks. It confirms that static variables, such as MGR and tax seniority, remain structuring axes for grouping. However, the rigidity of its grid mechanically constrains its maximum compression capabilities.
    \item \textbf{Industrial optimum via Machine Learning:} To push the limits of compression, unsupervised algorithms stand out. While \textit{HDBSCAN} shows remarkable theoretical precision thanks to its fine management of local density, the nature of its hyperparameters makes targeting a precise volume particularly complex. The \textit{K-Means} algorithm thus emerges as the best operational compromise: it ensures direct control over the number of generated \textit{Model Points} while capturing the deformation of capital requirements with high fidelity.
    \item \textbf{Limits of dynamic approaches:} The \textit{Cash-Flow Matching} method proves efficient on central projections, but its smoothing of individual characteristics penalizes it during the evaluation of regulatory shocks. This deviation affects not only lapse risk but mechanically extends to all risk modules sensitive to policyholder behavior.
\end{itemize}

The scope of these conclusions remains subject to the study's assumptions. On one hand, the use of a stochastic generator, though based on empirical data, necessarily introduces a simplification compared to the complexity of a real portfolio. On the other hand, sensitivity tests were limited only to interest rate risk. An evaluation integrating other changes in projection assumptions would allow for a broader testing of aggregation methods. Furthermore, enriching the ALM model by implementing a more sophisticated profit-sharing algorithm and introducing a dynamic lapse law would constitute a particularly relevant research axis to test the robustness of groupings against highly non-linear behaviors.

This work opens several avenues for improvement in ALM modeling. The mechanics of the \textit{HDBSCAN} algorithm offer a particularly relevant research path. Rather than imposing a defined number of \textit{Model Points}, this algorithm allows for the deduction of optimal volume directly from data density. Unsupervised learning thus allows for exceeding compression targets in favor of modeling driven by the very nature of financial liabilities. On the other hand, dynamic approaches could be enriched and calibrated on ALM projections to better capture Asset-Liability interactions and thus limit aggregation biases during regulatory shock evaluations. Finally, validating these methods on real portfolios would confirm the operational relevance of these aggregation techniques in an industrial context.